The mapper algorithm is a popular tool in topological data analysis (TDA) and 
visualization (TDV). Mapper takes as input a point cloud $X$, a clustering 
scheme $\text{Clust}: S \to Part(S)$, a filter (also called lens) function 
$f: X \to B$, and a cover of the filtered space $B$. The cover is used to 
break down the data into overlapping regions of our filtered data which, 
when pulled back by $f^{-1}$ allows for clustering within the overlapping 
regions in the original space. Mapper takes these clusters and returns the 
simplicial complex that is the nerve of the 'open cover' of these clusters.

This description of the mapper algorithm is complex for those without a 
mathematical background and is unmotivated without knowledge of the nerve 
theorem. To address this the ShinyMappeR project provides an interactive
shinyapp tool for teaching, learning, and exploring the effect paramter 
choice has on the output. 

ShinyMappeR is written in the R programming language and uses the recent 
mappeR implementation of the mapper algorithm. There are many default 
datasets (e.g., circles, figure-eights, spirals, and barbells) lens 
functions (coordinate projections, eccentricity, PCA-1), and clustering 
methods (DBSCAN, Agglomerative Clustering Methods: Single Linkage, Average Linkage,
Complete Linkage, and Ward's Linkage).

The structure of the application is simple, containing two sections: a user 
interface which to select the dataset and parameter choices for mapper, 
along with a visualization window which displays the original dataset, lens function,  
preimages of each open cover of the filtered space, and an identifaction of nodes
from the mapper complex to clusters within the pullback of each open cover.

ShinyMappeR also contains a visualization for the recent TDA mapper method G-Mapper.
[Desc. of G-Mapper] [Desc. of visualization].
